The electroweak (EWK) theory provides a unified description of the electromagnetic and weak interactions, obtained by imposing local gauge invariance under the group $SU(2)_L \times U(1)_Y$. 

The weak interaction couples universally to all fermions, connects different quark flavors through the Cabibbo-Kobayashi-Maskawa (CKM) mixing matrix, and is the only interaction observed to violate charge (C), parity (P), and combined charge-parity (CP) symmetries. The violation of parity~\cite{Lee:1956qn,Wu:1957my} manifests itself in the fact that the weak interaction acts differently on fermions depending on their \textit{chirality}, defined through the matrix $\gamma^5 = i\gamma^0\gamma^1\gamma^2\gamma^3$ built from the Dirac matrices of Equation \ref{eq:dirac_lagrangian}. Any fermion field can be decomposed into a left-handed and a right-handed component, obtained with the projection operators $P_L = (1-\gamma^5)/2$ and $P_R = (1+\gamma^5)/2$, and only the left-handed components of the fermions, or the right-handed components of the antifermions, transform under $SU(2)_L$, as indicated by its subscript $L$. In the limit of a massless particle, chirality coincides with helicity, the projection of the spin onto the direction of motion.

Accordingly, the fermion fields of each generation are arranged into a left-handed doublet, which transforms under $SU(2)_L$, and right-handed singlets, which do not — two for quarks, and only one for leptons, since the SM contains no right-handed neutrino —,
\begin{equation}
    \Psi_L \equiv P_L\begin{pmatrix}\psi \\ \psi'\end{pmatrix}=\begin{pmatrix}\psi_L \\ \psi'_L\end{pmatrix}, \qquad
    \psi_R \equiv P_R\,\psi, \qquad
    \psi'_R \equiv P_R\,\psi',
    \label{fermionic_fields}
\end{equation}
%
where $\psi$ and $\psi'$ denote the two members of a doublet, either the neutrino and the charged lepton or the up- and down-type quarks. The charge associated with $SU(2)_L$ is the weak isospin $T_3$, which takes the values $T_3=+1/2$ for the upper and $T_3=-1/2$ for the lower component of the doublets, and $T_3=0$ for the right-handed singlets. As in QCD, local invariance requires one gauge field for each generator of the group, so that three fields $W^i_\mu$, with $i=1,2,3$, are introduced, whose field-strength tensors $W^i_{\mu\nu}$ contain, like those of the gluons, a quadratic term that gives rise to cubic and quartic self-interactions. However, the neutral field $W^3_\mu$ couples only to left-handed fermions and therefore cannot be identified with the photon, which couples equally to both chiralities. The electromagnetic interaction is therefore incorporated into the same symmetry by introducing an additional abelian group, $U(1)_Y$, whose charge is the weak hypercharge $Y$, carried by both chiralities and related to the electric charge through the Gell-Mann-Nishijima relation, $Q = T_3 + Y/2$. Local invariance under the full group then requires the covariant derivative
\begin{equation}
    D_\mu = \partial_\mu + i g\, T_i\, W^i_\mu + i g'\, \frac{Y}{2}\, B_\mu,
    \label{eq:ew_covariant_derivative}
\end{equation}
%
where $g$ is the weak coupling constant, $T_i = \sigma_i/2$, with $\sigma_i$ the Pauli matrices, are the generators of $SU(2)_L$, which vanish for the right-handed singlets, and $B_\mu$ is the gauge field of $U(1)_Y$, with coupling constant $g'$ and the abelian field-strength tensor $B_{\mu\nu} = \partial_\mu B_\nu - \partial_\nu B_\mu$. The EWK Lagrangian finally reads
\begin{equation}
    \mathcal{L}_{\text{EWK}} =
    \underbrace{\bar{\Psi}_L\, i \slashed{D}\, \Psi_L}_{\substack{\text{left-handed}\\ \text{doublets}}}
    + \underbrace{\bar{\psi}_R\, i \slashed{D}\, \psi_R + \bar{\psi}'_R\, i \slashed{D}\, \psi'_R}_{\substack{\text{right-handed}\\ \text{singlets}}}
    \underbrace{{}- \frac{1}{4}\, W^i_{\mu\nu} W_i^{\mu\nu} - \frac{1}{4}\, B_{\mu\nu} B^{\mu\nu}}_{\substack{\text{gauge}\\ \text{kinetic terms}}}.
    \label{EW_Lagrangian}
\end{equation}
%
The fermion terms contain, through the covariant derivative, both the free propagation of the fermions and their interactions with the gauge fields. These fields, however, do not correspond directly to the observed bosons, which arise as combinations of them. The first two $SU(2)_L$ fields combine into the charged fields $W^\pm_\mu = (W^1_\mu \mp i\,W^2_\mu)/\sqrt{2}$, associated with the $W^\pm$ bosons, while the two neutral fields, $W^3_\mu$ and $B_\mu$, mix into the fields of the $Z$ boson and the photon through a rotation by the Weinberg angle $\theta_W$, a free parameter of the SM,
\begin{equation}
    \begin{pmatrix} A_\mu \\ Z_\mu \end{pmatrix} = \begin{pmatrix} \cos\theta_W & \sin\theta_W \\ -\sin\theta_W & \cos\theta_W \end{pmatrix} \begin{pmatrix} B_\mu \\ W^3_\mu \end{pmatrix}, \qquad \tan\theta_W = \frac{g'}{g}.
    \label{eq:weinberg_mixing}
\end{equation}
%
Written in terms of these fields, the EWK Lagrangian describes two types of interactions between fermions and gauge bosons. The charged currents, mediated by the $W^\pm$ bosons, act only on the left-handed components and change the flavor of the fermion within the doublet. The neutral currents are mediated by the photon, which couples to the electric charge with a strength $e = g\sin\theta_W = g'\cos\theta_W$, identified as the unit of electric charge, and by the $Z$ boson, which couples to both chiralities, but with different strengths. The gauge kinetic terms, in turn, give rise to the self-interactions of the $W^\pm$ bosons with the photon and the $Z$ boson, such as the $WW\gamma$ and $WWZ$ vertices, while photons do not interact with each other.\medskip



l 191: es in the case of the gluones and the phogon -> as in the case of the gluons
which was allowed in QCD and QED --> which was allowed in QCD



STRONG INTERACTIN: l73 to 110

The strong interaction is described by QCD, a gauge theory based on the non-abelian group $SU(3)_C$, where the subscript $C$ refers to color, the conserved charge associated with the symmetry. Its construction starts from the Lagrangian density of a free fermion field, the Dirac Lagrangian,
\begin{equation}
    \mathcal{L}_{\text{Dirac}} = \bar{\psi}(x)\,(i\slashed{\partial} - m)\,\psi(x),
    \label{eq:dirac_lagrangian}
\end{equation}
%
where $\psi(x)$ is the quark spinor field at the spacetime coordinates $x$, $m$ is the quark mass, and $\slashed{\partial} = \gamma^\mu \partial_\mu$ denotes the contraction of the Dirac matrices $\gamma^\mu$ with the spacetime derivatives. This Lagrangian is invariant under global $SU(3)_C$ transformations of the quark field, but not under local ones, whose parameters vary from one spacetime point to another. Local gauge invariance is restored by replacing $\partial_\mu$ with the covariant derivative $D_\mu = \partial_\mu + i g_s\, (\lambda_a/2)\, G^a_\mu$, where $\lambda_a$ are the eight Gell-Mann matrices that generate the group and $g_s$ is the strong coupling constant. This introduces one vector field $G^a_\mu$ for each generator, the eight gluon fields, and the QCD Lagrangian reads
\begin{equation}
    \mathcal{L}_{\text{QCD}} =
    \underbrace{\bar{\psi}\,(i\slashed{\partial} - m)\,\psi}_{\substack{\text{free quark}\\ \text{propagation}}}
    \underbrace{{}-g_s\, \bar{\psi}\gamma^\mu \frac{\lambda_a}{2}\psi\, G^a_\mu}_{\substack{\text{quark-gluon}\\ \text{interaction}}}
    \underbrace{{}-\frac{1}{4}\, G^a_{\mu\nu} G_a^{\mu\nu}}_{\substack{\text{gluon kinetic}\\ \text{term}}}.
    \label{QCD_Lagrangian}
\end{equation}
%
The first term is the Dirac Lagrangian of Equation \ref{eq:dirac_lagrangian}, while the second, which arises from the covariant derivative, describes the interaction between quarks and gluons, with a strength usually expressed through $\alpha_s = g_s^2/4\pi$. The third is the kinetic term of the gluons, built from the field-strength tensor $G^a_{\mu\nu} = \partial_\mu G^a_\nu - \partial_\nu G^a_\mu - g_s f^{abc}\, G^b_\mu G^c_\nu$, where $f^{abc}$ are the structure constants of $SU(3)_C$. Its last term, a direct consequence of the non-abelian nature of the group, gives rise to cubic and quartic gluon self-interactions. In contrast, a gluon mass term, of the form $m_g^2\, G^a_\mu G_a^\mu$, would break local gauge invariance, so that gluons are necessarily massless.


Thus, local gauge invariance not only requires the existence of the gluon fields, but also determines their interactions with quarks and among themselves.



\subsubsection{Electroweak interaction}
%
The electroweak theory provides a unified description of the electromagnetic and weak interactions, built on the same principle as QCD, the requirement of local gauge invariance, but under the group $SU(2)_L \times U(1)_Y$. However, its construction its more complex because the group is not fixed by a single charge, as color is in QCD, but by the observed properties of the weak interaction. The $SU(2)_L$ factor reflects the fact that the weak interaction acts only on the left-handed components of the fermions, which it connects in pairs, while the additional $U(1)_Y$ factor is required because the weak interaction cannot be described independently of electromagnetism.

The weak interaction couples universally to all fermions, connects different quark flavors through the Cabibbo-Kobayashi-Maskawa (CKM) mixing matrix, and is the only interaction observed to violate charge (C), parity (P), and combined charge-parity (CP) symmetries. The violation of parity~\cite{Lee:1956qn,Wu:1957my} reflects the fact that it acts differently on fermions depending on their \textit{chirality}, a property defined through the matrix $\gamma^5 = i\gamma^0\gamma^1\gamma^2\gamma^3$, built from the Dirac matrices of Equation \ref{eq:dirac_lagrangian}. Any fermion field can be decomposed into a left-handed and a right-handed component, obtained with the projection operators $P_L = (1-\gamma^5)/2$ and $P_R = (1+\gamma^5)/2$, and in the limit of a massless particle chirality coincides with helicity, the projection of the spin onto the direction of motion. Experimentally, only the left-handed components of the fermions — or the right-handed components of the antifermions — transform under $SU(2)_L$, as indicated by its subscript $L$. 

Weak processes proceed through both charged currents, which change the electric charge of the fermion, and neutral currents, which leave it unchanged and, as shown below, arise as a consequence of the symmetry. In the charged currents, each fermion is transformed into a partner of the same generation, so that the two partners behave as two states of the same object, which the $SU(2)_L$ transformations rotate into one another. Accordingly, the fermion fields of each generation are arranged into a left-handed doublet, which transforms under $SU(2)_L$, and right-handed singlets, which do not — two for quarks, and only one for leptons, since the SM contains no right-handed neutrino:


These charged interactions Y QUE PASA CON LAS NO CHARGED? transform each fermion into a partner of the same generation, as in the $\beta$ decay of the neutron, HAY QUE MENCIONAR EL BETA DECAY ? in which a $d$ quark is transformed into a $u$ quark while an electron and an antineutrino are produced. The two partners therefore behave as two states of the same object, which the $SU(2)_L$ transformations rotate into one another. ME GUSTA LA INTRODUCCION ANTES DEL ACCORDINGLY, PERO NO EXPLICAR EL BETA DECAY. Accordingly, the fermion fields of each generation are arranged into a left-handed doublet, which transforms under $SU(2)_L$, and right-handed singlets, which do not — two for quarks, and only one for leptons, since the SM contains no right-handed neutrino:
\begin{equation}
    \Psi_L \equiv P_L\begin{pmatrix}\psi \\ \psi'\end{pmatrix}=\begin{pmatrix}\psi_L \\ \psi'_L\end{pmatrix}, \qquad
    \psi_R \equiv P_R\,\psi, \qquad
    \psi'_R \equiv P_R\,\psi',
    \label{fermionic_fields}
\end{equation}
%
where $\psi$ and $\psi'$ denote the two members of a doublet, either the neutrino and the charged lepton or the up- and down-type quarks. Leptons and quarks thus belong to separate doublets, and since no interaction of the SM connects them, a quark can never be transformed into a lepton or vice versa. The charge associated with $SU(2)_L$ is the weak isospin $T_3$, which distinguishes the two members of the doublet, with $T_3=+1/2$ for the upper and $T_3=-1/2$ for the lower component, and is zero for the right-handed singlets.

As in QCD, local invariance under $SU(2)_L$ requires one gauge field for each generator of the group, so that three fields $W^i_\mu$, with $i=1,2,3$, are introduced, with a common coupling constant $g$. The generators, $T_i = \sigma_i/2$, with $\sigma_i$ the Pauli matrices, act on the doublets in two different ways. $T_1$ and $T_2$ exchange the two members of the doublet, so that $W^1_\mu$ and $W^2_\mu$ mediate the charged interactions, while $T_3$ leaves each member unchanged, so that $W^3_\mu$ mediates a neutral interaction. Moreover, since $SU(2)_L$ is non-abelian, the field-strength tensors $W^i_{\mu\nu}$ contain, like that of the gluons, a term quadratic in the fields, which gives rise to self-interactions among the $W^i_\mu$. However, $SU(2)_L$ alone cannot describe the observed interactions, since its neutral field $W^3_\mu$ couples only to left-handed fermions, in proportion to $T_3$, and therefore cannot be identified with the photon, which couples to the electric charge $Q$ of both chiralities. AQUI YA QUE NO SE PONE LA ECUACION, QUIZAS SE PUEDA REFERENCIAR SI HAY UNA ECUACION ANALOGA EN STRONG INTERACTION? Y ME GUSSTABA LA PARTE DE DECIR QUE  ENTONCES LA INTERACCION ELECTROMAGNETICA TENIA QUE SER INCORPORADA EN LA MISMA SIMETRÍA\medskip

The electromagnetic interaction is therefore incorporated into the same symmetry through an additional abelian group, $U(1)_Y$, with a gauge field $B_\mu$ and a coupling constant $g'$. Its charge, the weak hypercharge $Y$, is carried by both chiralities and is assigned to each multiplet so as to reproduce the observed electric charges through the Gell-Mann-Nishijima relation, $Q = T_3 + Y/2$. The electric charge is thus not an independent quantity but a combination of the charges of $SU(2)_L$ and $U(1)_Y$, so that electromagnetism is contained in the electroweak symmetry. Local invariance under the full group then requires the covariant derivative
\begin{equation}
    D_\mu = \partial_\mu + i g\, T_i\, W^i_\mu + i g'\, \frac{Y}{2}\, B_\mu,
    \label{eq:ew_covariant_derivative}
\end{equation}
%
where $T_i = 0$ for the right-handed singlets, which therefore couple only to $B_\mu$, whose field-strength tensor, $B_{\mu\nu} = \partial_\mu B_\nu - \partial_\nu B_\mu$, contains no quadratic term, since $U(1)_Y$ is abelian. Combining all these ingredients, the EWK Lagrangian reads Y TAMPOCO DICES NADA DE LA HYPERCHARGE FIELD
\begin{equation}
    \mathcal{L}_{\text{EWK}} =
    \underbrace{\bar{\Psi}_L\, i \slashed{D}\, \Psi_L}_{\substack{\text{left-handed}\\ \text{doublets}}}
    + \underbrace{\bar{\psi}_R\, i \slashed{D}\, \psi_R + \bar{\psi}'_R\, i \slashed{D}\, \psi'_R}_{\substack{\text{right-handed}\\ \text{singlets}}}
    \underbrace{{}- \frac{1}{4}\, W^i_{\mu\nu} W_i^{\mu\nu} - \frac{1}{4}\, B_{\mu\nu} B^{\mu\nu}}_{\substack{\text{gauge}\\ \text{kinetic terms}}},
    \label{EW_Lagrangian}
\end{equation}
%
which has the same structure as the QCD Lagrangian of Equation \ref{QCD_Lagrangian}. The fermion terms contain, through the covariant derivative, both the free propagation of the fermions and their interactions with the gauge fields, while the kinetic terms describe the propagation of the gauge fields and, through the quadratic term of $W^i_{\mu\nu}$, their self-interactions. Unlike in QCD, however, it contains no mass terms, for reasons discussed below.

The fields $W^i_\mu$ and $B_\mu$ are those dictated by the symmetry, but they do not correspond directly to the observed bosons, which arise as combinations of them. The first two $SU(2)_L$ fields combine into $W^\pm_\mu = (W^1_\mu \mp i\,W^2_\mu)/\sqrt{2}$, which transform the lower member of a doublet into the upper one and vice versa, changing $T_3$ by one unit. Since both members of a doublet have the same hypercharge, the relation $Q = T_3 + Y/2$ implies that they also change the electric charge by one unit, so these fields describe the charged $W^\pm$ bosons. The two neutral fields, $W^3_\mu$ and $B_\mu$, couple to $T_3$ and to $Y$, respectively, so neither of them couples to the electric charge alone. The photon is instead the combination of the two that couples to $Q = T_3 + Y/2$, and the $Z$ boson the orthogonal one, both obtained through a rotation by the Weinberg angle $\theta_W$,
\begin{equation}
    \begin{pmatrix} A_\mu \\ Z_\mu \end{pmatrix} = \begin{pmatrix} \cos\theta_W & \sin\theta_W \\ -\sin\theta_W & \cos\theta_W \end{pmatrix} \begin{pmatrix} B_\mu \\ W^3_\mu \end{pmatrix}, \qquad \tan\theta_W = \frac{g'}{g},
    \label{eq:weinberg_mixing}
\end{equation}
%
where the requirement that the photon couple to $T_3$ and to $Y/2$ with the same strength fixes the angle through the ratio of the two coupling constants, which are free parameters of the SM. Rewriting the fermion terms of Equation \ref{EW_Lagrangian} in terms of these fields yields the interactions of the fermions with the physical gauge bosons,
\begin{equation}
    \begin{split}
    \mathcal{L}_{\text{CC}} &= -\frac{g}{\sqrt{2}} \left(\bar{\psi}_L\gamma^\mu\psi'_L\, W^+_\mu + \bar{\psi}'_L\gamma^\mu\psi_L\, W^-_\mu\right), \\
    \mathcal{L}_{\text{NC}} &= -e\,Q\,\bar{\psi}\gamma^\mu\psi\, A_\mu - \frac{g}{\cos\theta_W}\,\bar{\psi}\gamma^\mu\left(T_3\, P_L - Q\sin^2\theta_W\right)\psi\, Z_\mu,
    \end{split}
    \label{eq:ew_currents}
\end{equation}
%
Y QUE PASA AQUI CON FREE Y KIN?
where a sum over all fermions is implied. The charged currents, mediated by the $W^\pm$ bosons, act only on the left-handed components and change the flavor of the fermion within the doublet, as required by the observations from which the construction started. The neutral currents contain two terms. The first is the interaction of the photon with the electric charge, with a strength $e = g\sin\theta_W = g'\cos\theta_W$ that is identified as the unit of electric charge, and reproduces quantum electrodynamics (QED), which thus emerges as a part of the electroweak theory. The second describes the $Z$ boson, which, being a combination of $W^3_\mu$ and $B_\mu$, couples to both chiralities, but with different strengths, since its $T_3$ term acts only on the left-handed components. Finally, the gauge kinetic terms, written in terms of the physical fields, give rise to trilinear and quartic self-interactions that always involve a pair of $W$ bosons, such as the $WW\gamma$ and $WWZ$ vertices, while photons do not interact with each other.\medskip


